\section{Week 4: The Transformer — Attention as the Core Computational Primitive}

\subsection{Learning Objectives}

By the end of this week, you should be able to:
\begin{itemize}
  \item Derive scaled dot-product self-attention
  \item Explain the roles of queries, keys, and values
  \item Understand the structure of a Transformer block
  \item Reason about residual streams and layer normalization
  \item Explain why Transformers scale better than prior sequence models
\end{itemize}

\subsection{LLM Components Covered}

\textbf{Core architecture}
\begin{itemize}
  \item Token embeddings + positional information
  \item Self-attention (Q, K, V)
  \item Multi-head attention
  \item Feed-forward (MLP) blocks
  \item Residual connections and layer normalization
\end{itemize}

\textbf{Systems relevance}
\begin{itemize}
  \item Parallelism and scalability
  \item Memory and compute tradeoffs
  \item Attention as content-addressable memory
\end{itemize}

\subsection{Conceptual Core: Why Attention Replaced Recurrence}

Prior sequence models (RNNs, LSTMs) process tokens sequentially, limiting:
\begin{itemize}
  \item parallelism,
  \item long-range dependency modeling,
  \item scalability.
\end{itemize}

Transformers eliminate recurrence entirely.
Instead, they allow each token to attend directly to every other token in the context window.

\subsection{Scaled Dot-Product Self-Attention}

Given an input sequence represented as vectors
\[
X \in \mathbb{R}^{T \times d}
\]

we compute:

\[
Q = XW_Q,\quad K = XW_K,\quad V = XW_V
\]

Attention weights are:

\[
A = \mathrm{softmax}\left(\frac{QK^\top}{\sqrt{d_k}}\right)
\]

and outputs:

\[
\mathrm{Attention}(Q,K,V) = AV
\]

Interpretation:
\begin{itemize}
  \item Queries ask questions
  \item Keys advertise content
  \item Values carry information
\end{itemize}

This is best understood as \emph{learned, differentiable lookup}.

\subsection{Multi-Head Attention}

Instead of a single attention operation, Transformers use multiple heads:

\[
\mathrm{MHA}(X) = \mathrm{Concat}(h_1,\dots,h_H)W_O
\]

Each head can learn a different relational pattern:
\begin{itemize}
  \item syntax,
  \item coreference,
  \item positional relationships.
\end{itemize}

\subsection{The Transformer Block}

A decoder-style Transformer block consists of:

\begin{enumerate}
  \item LayerNorm
  \item Masked self-attention
  \item Residual connection
  \item LayerNorm
  \item Feed-forward network (MLP)
  \item Residual connection
\end{enumerate}

The residual stream allows information to flow across layers with minimal degradation.

\subsection{Why Positional Information Is Required}

Attention is permutation-invariant.
Without positional encoding, token order is lost.

Position information is injected via:
\begin{itemize}
  \item sinusoidal encodings,
  \item learned embeddings,
  \item rotary position encodings (conceptually).
\end{itemize}

\subsection{Systems Engineering Implications}

\textbf{Requirements}
\begin{itemize}
  \item Efficient parallel computation
  \item Stable training across depth
\end{itemize}

\textbf{Architecture}
\begin{itemize}
  \item Attention scales as $O(T^2)$
  \item KV-cache reduces inference cost
\end{itemize}

\textbf{Verification \& Validation}
\begin{itemize}
  \item Inspect attention patterns
  \item Check masking correctness
\end{itemize}

\textbf{Operations}
\begin{itemize}
  \item Context window management
  \item Memory and latency monitoring
\end{itemize}

\subsection{Human Factors and Failure Modes}

Common misunderstandings:
\begin{itemize}
  \item ``Attention equals explanation'' (it does not)
  \item Overinterpreting attention visualizations
\end{itemize}

Attention weights reflect internal computation, not human-interpretable reasoning.

\subsection{Key References}

\begin{itemize}
  \item \citet{vaswani2017attention} — Transformer architecture
\end{itemize}

\subsection{Recommended University Lectures}

\begin{itemize}
  \item Stanford CS224N — Transformers and attention  
  \url{https://web.stanford.edu/class/cs224n/}

  \item MIT 6.S191 — Attention mechanisms  
  \url{https://introtodeeplearning.mit.edu}
\end{itemize}

\subsection{Practitioner Lab}

See \texttt{labs/week04/week04\_lab\_transformer\_blocks.ipynb}.

\subsection{Week 4 Takeaway}

\begin{quote}
The Transformer turns sequence modeling into parallel, differentiable memory access.
\end{quote}

